% AA-68
\chapter[Charaktery reprezentacji]{Charaktery formalne reprezentacji półprostych algebr Liego}
\label{aa:ch:lie-formal-characters}

Klasyfikacja z AA-67 opisuje reprezentacje nieprzywiedlne przez ich wagi najwyższe, ale nie podaje jeszcze krotności pozostałych wag. Zapiszemy wszystkie krotności jako współczynniki jednego elementu pierścienia. Pozwoli to również badać rozkład iloczynów tensorowych.

W całym rozdziale $K$ jest algebraicznie domkniętym ciałem charakterystyki zero, $\mathfrak g$ niezerową skończenie wymiarową półprostą algebrą Liego nad $K$, $\mathfrak h$ maksymalną podalgebrą toralną, $\Delta=(\alpha_1,\ldots,\alpha_r)$ bazą układu korzeni, $W$ grupą Weyla, a $P,P_+,\preceq$ mają znaczenie z Definicji~\ref{aa:def:lie-dominant-integral}. Przez $L(\lambda)$ oznaczamy moduł nieprzywiedlny o wadze najwyższej $\lambda\in P_+$ z Twierdzenia~\ref{aa:tw:lie-highest-classification}.

\section{Pierścień grupowy kraty wag}

\begin{definicja}
\label{aa:def:lie-formal-character}
Dla addytywnej kraty wag $P$ pierścień grupowy $\mathbb Z[P]$ jest grupą abelową swobodną o bazie symboli $e^\mu$, $\mu\in P$, z mnożeniem określonym przez $e^\mu e^\nu=e^{\mu+\nu}$ i rozszerzonym przez rozdzielność. Dla skończenie wymiarowego $\mathfrak g$-modułu $V$ definiujemy jego charakter formalny
\[
 \operatorname{ch} V=\sum_{\mu\in P}(\dim_K V_\mu)e^\mu\in\mathbb Z[P].
\]
Suma ma skończony nośnik na mocy Twierdzenia~\ref{aa:tw:lie-weight-decomposition}. Symbole $e^\mu$ są formalne, a nie wartościami funkcji wykładniczej.
\end{definicja}

Jeżeli $\omega_1,\ldots,\omega_r$ są wagami fundamentalnymi, to $e^{\sum_i n_i\omega_i}\mapsto x_1^{n_1}\cdots x_r^{n_r}$ daje izomorfizm $\mathbb Z[P]\cong\mathbb Z[x_1^{\pm1},\ldots,x_r^{\pm1}]$: jednoznaczność współrzędnych w bazie wag fundamentalnych daje bijekcję baz, a dodawanie wykładników jest zgodne z mnożeniem. W szczególności nie należy utożsamiać symbolu $e^0=1$ z wektorem wagi zero.

\begin{twierdzenie}
\label{aa:tw:lie-character-operations}
Niech $\mathfrak g$ będzie niezerową skończenie wymiarową półprostą algebrą Liego nad algebraicznie domkniętym ciałem $K$ charakterystyki zero. Ustalmy maksymalną podalgebrę toralną i bazę korzeni prostych, a przez $P,W$ oznaczmy odpowiadające im kratę wag i grupę Weyla. Dla skończenie wymiarowych $\mathfrak g$-modułów $U,V$ zachodzą równości
\begin{align*}
 \operatorname{ch}(U\oplus V)&=\operatorname{ch}U+\operatorname{ch}V,\\
 \operatorname{ch}(U\otimes_K V)&=\operatorname{ch}U\,\operatorname{ch}V,\\
 \operatorname{ch}(V^*)&=\sum_{\mu\in P}(\dim_K V_\mu)e^{-\mu}.
\end{align*}
Ponadto $\operatorname{ch}V$ należy do pierścienia niezmienników $\mathbb Z[P]^W$, gdzie $w(e^\mu)=e^{w\mu}$, a homomorfizm $\varepsilon\colon\mathbb Z[P]\to\mathbb Z$ określony przez $\varepsilon(e^\mu)=1$ spełnia $\varepsilon(\operatorname{ch}V)=\dim_K V$.
\end{twierdzenie}
\begin{proof}
Rozkłady wagowe $U$ i $V$ dają $(U\oplus V)_\gamma=U_\gamma\oplus V_\gamma$. Dla iloczynu tensorowego działanie jest określone wzorem $t(u\otimes v)=tu\otimes v+u\otimes tv$ dla $t\in\mathfrak g$. Jeżeli $u\in U_\mu$, $v\in V_\nu$ oraz $t\in\mathfrak h$, to $t(u\otimes v)=(\mu+\nu)(t)u\otimes v$. Bazy przestrzeni wagowych obu czynników dają bazę iloczynu tensorowego złożoną z takich wektorów. Zatem
\[
 (U\otimes_K V)_\gamma=
 \bigoplus_{\mu+\nu=\gamma} U_\mu\otimes_K V_\nu.
\]
Rzeczywiście, grupując elementy tej bazy według sumy wag, otrzymujemy sumę prostą; porównanie składowych równania $tz=\gamma(t)z$ dla wszystkich $t\in\mathfrak h$ wyklucza każdą składową o innej wadze. Wymiary składowych są iloczynami wymiarów czynników, co daje współczynniki iloczynu charakterów.

Na przestrzeni sprzężonej $V^*$ działanie określamy przez $(t\varphi)(v)=-\varphi(tv)$. Funkcjonał na $V_\mu$, przedłużony przez zero na pozostałych składowych, ma wagę $-\mu$. Ponieważ rozkład $V$ jest skończony, takie funkcjonały rozpinają $V^*$ i dają podany wzór. Każde odbicie proste zachowuje kratę $P$, gdyż $s_i\mu=\mu-\mu(h_i)\alpha_i$ i korzenie należą do $P$. Stąd działanie $W$ permutuje bazę pierścienia i zachowuje mnożenie. Twierdzenie~\ref{aa:tw:lie-weight-reflection} daje $\dim V_{w\mu}=\dim V_\mu$, więc charakter jest niezmienniczy. Wreszcie $\varepsilon$ zachowuje mnożenie na elementach bazy, a więc na całym pierścieniu; suma wymiarów przestrzeni wagowych jest wymiarem $V$.
\end{proof}

\section{Trójkątność i pierścień niezmienników}

\begin{lemat}
\label{aa:lem:lie-character-triangular}
Niech $\mathfrak g$ będzie niezerową skończenie wymiarową półprostą algebrą Liego nad algebraicznie domkniętym ciałem charakterystyki zero. Ustalmy maksymalną podalgebrę toralną i bazę korzeni prostych. Dla każdej $\lambda\in P_+$ charakter modułu $L(\lambda)$ ma postać
\[
 \operatorname{ch}L(\lambda)=e^\lambda+
 \sum_{\mu\prec\lambda}m_{\lambda,\mu}e^\mu,
 \qquad m_{\lambda,\mu}\in\mathbb Z_{\ge0},
\]
gdzie $\mu\prec\lambda$ oznacza $\mu\preceq\lambda$ i $\mu\ne\lambda$. Charakter zawiera tylko skończenie wiele składników niezerowych.
\end{lemat}
\begin{proof}
Z Twierdzenia~\ref{aa:tw:verma-universal} moduł Verma $M(\lambda)$ ma jednowymiarową najwyższą przestrzeń wagową i wszystkie jego wagi są $\preceq\lambda$. Z Lematu~\ref{aa:lem:verma-simple-quotient} obraz generatora w $L(\lambda)$ jest niezerowy. Rozkład wagowy ilorazu (argument z Lematu~\ref{aa:lem:verma-integrable}) zachowuje więc jednowymiarową składową najwyższą i nie wprowadza nowych wag. Skończoność wymiaru $L(\lambda)$ wynika z Twierdzenia~\ref{aa:tw:lie-highest-classification}. Współczynniki są wymiarami przestrzeni wagowych, co dowodzi tezy.
\end{proof}

\begin{twierdzenie}
\label{aa:tw:lie-character-basis}
Niech $\mathfrak g$ będzie niezerową skończenie wymiarową półprostą algebrą Liego nad algebraicznie domkniętym ciałem charakterystyki zero. Ustalmy maksymalną podalgebrę toralną i bazę korzeni prostych; $P$, $P_+$ i $W$ oznaczają kratę wag, wagi całkowite dominujące i grupę Weyla. Rodzina $\{\operatorname{ch}L(\lambda):\lambda\in P_+\}$ jest bazą grupy abelowej $\mathbb Z[P]^W$.
\end{twierdzenie}
\begin{proof}
Niech $\ell$ będzie funkcjonałem rzeczywistym na $P\otimes_{\mathbb Z}\mathbb R$ określonym przez $\ell(\alpha_i)=1$. Korzenie proste tworzą bazę tej przestrzeni: macierz ich współrzędnych w bazie wag fundamentalnych jest całkowita i odwracalna. Dla $\nu\prec\lambda$ mamy $\ell(\nu)\le\ell(\lambda)-1$. Dla $\lambda\in P_+$ zachodzi $\ell(\lambda)\ge0$. Istotnie, dowód Twierdzenia~\ref{aa:tw:lie-highest-classification} wykazał, że $|W|\lambda=\sum_i b_i\alpha_i$ z $b_i\in\mathbb Z_{\ge0}$; zatem $\ell(\lambda)=\sum_i b_i/|W|$. Ponadto $\ell(P)\subset D^{-1}\mathbb Z$ dla pewnego dodatniego całkowitego $D$, gdyż odwrotność macierzy całkowitej o niezerowym wyznaczniku ma współczynniki wymierne.

Rozważmy niezerowy element $F=\sum_\mu c_\mu e^\mu$ pierścienia niezmienników. Wybierzmy wagę $\lambda$ o największej wartości $\ell$ w skończonym nośniku $F$. Jest ona dominująca: gdyby $\lambda(h_i)<0$, to $s_i\lambda=\lambda-\lambda(h_i)\alpha_i$ miałaby większą wartość $\ell$ i należałaby do nośnika, ponieważ $F$ jest niezmienniczy. Otrzymujemy sprzeczność. Zatem $\lambda\in P_+$.

Odejmijmy $c_\lambda\operatorname{ch}L(\lambda)$. Wynik nadal jest niezmienniczy, współczynnik przy $e^\lambda$ znika, a wszystkie nowe wagi mają wartość $\ell$ mniejszą od $\ell(\lambda)$ na mocy Lematu~\ref{aa:lem:lie-character-triangular}. Dla każdej wartości maksymalnej występuje skończenie wiele wag (nośnik w danym kroku jest skończony), a żaden krok na tej wysokości nie wprowadza nowych wag tej samej wysokości. Po skończenie wielu krokach na danej wysokości albo kończymy, albo maksymalna wysokość maleje. Każda niezerowa reszta ma maksymalną wysokość nieujemną, ponieważ jej waga maksymalna jest dominująca. Możliwych wysokości między zerem a początkowym maksimum jest skończenie wiele, gdyż należą do $D^{-1}\mathbb Z$. Procedura kończy się więc zerem i daje skończoną kombinację całkowitą charakterów.

Dla niezależności liniowej załóżmy, że skończona kombinacja $\sum_{\lambda} a_\lambda\operatorname{ch}L(\lambda)$ jest zerowa i nie wszystkie $a_\lambda$ są zerowe. Wybierzmy $\lambda$ o największej wysokości wśród wag z $a_\lambda\ne0$. Charakter o innej wadze najwyższej o tej samej lub mniejszej wysokości nie zawiera $e^\lambda$, ponieważ każda waga różna od najwyższej ma wysokość ściśle mniejszą. Współczynnik przy $e^\lambda$ wynosi zatem $a_\lambda\ne0$, sprzecznie z założeniem.
\end{proof}

\begin{wniosek}
\label{aa:cor:lie-character-determines}
Dla niezerowej skończenie wymiarowej półprostej algebry Liego nad algebraicznie domkniętym ciałem charakterystyki zero dwa skończenie wymiarowe moduły są izomorficzne wtedy i tylko wtedy, gdy mają równe charaktery formalne względem tej samej maksymalnej podalgebry toralnej.
\end{wniosek}
\begin{proof}
Izomorfizm zachowuje przestrzenie wagowe. Odwrotnie, na mocy Twierdzenia~\ref{aa:tw:lie-weyl-complete-reducibility} każdy z modułów jest skończoną sumą prostą modułów nieprzywiedlnych. Twierdzenie~\ref{aa:tw:lie-highest-classification} identyfikuje ich klasy z wagami dominującymi. Równość charakterów oraz niezależność z Twierdzenia~\ref{aa:tw:lie-character-basis} dają równe krotności każdej klasy w obu sumach. Dobierając izomorfizmy odpowiednich składników, otrzymujemy izomorfizm modułów. Dotyczy to również modułu zerowego.
\end{proof}

\section{Przykład i rozkład iloczynu tensorowego}

\begin{przyklad}[Rząd jeden]
\label{aa:prz:lie-character-sl2}
Dla $\mathfrak g=\mathfrak{sl}_2(K)$ oznaczmy przez $\omega$ wagę fundamentalną i przez $x=e^\omega$. Korzeń prosty to $2\omega$, a odbicie proste wysyła $x$ na $x^{-1}$. Z modelu AA-67 otrzymujemy
\[
 C_m=\operatorname{ch}L(m\omega)=x^m+x^{m-2}+\cdots+x^{-m}
 \qquad(m\ge0).
\]
Dla $a\ge b\ge0$ zachodzi $C_a C_b=\sum_{j=0}^{b}C_{a+b-2j}$. Istotnie, współczynnik przy $x^{a+b-2k}$ w iloczynie jest liczbą par $(i,j)$ takich, że $0\le i\le a$, $0\le j\le b$, $i+j=k$. Dla $0\le k\le a+b$ jest to
$\min(b,k)-\max(0,k-a)+1$. W sumie po prawej stronie wykładnik ten występuje raz w składniku o indeksie $j$ dokładnie wtedy, gdy $0\le j\le b$ oraz $j\le k\le a+b-j$. Liczba takich indeksów wynosi $\min(b,k,a+b-k)+1$. Obie liczby są równe: dla $k\le a$ mamy $a+b-k\ge b$, a dla $k>a$ mamy $k>b$ i $a+b-k<b$. Poza zakresem $0\le k\le a+b$ oraz dla wykładników o innej parzystości oba współczynniki są zerowe. Wobec Twierdzenia~\ref{aa:tw:lie-character-operations} i Wniosku~\ref{aa:cor:lie-character-determines} dostajemy
\[
 L(a\omega)\otimes_K L(b\omega)
 \cong \bigoplus_{j=0}^b L((a+b-2j)\omega).
\]
Przypadek $a<b$ wynika z izomorfizmu $u\otimes v\mapsto v\otimes u$, zgodnego z działaniem algebry Liego.
\end{przyklad}

W wyższych rzędach trójkątna procedura z dowodu Twierdzenia~\ref{aa:tw:lie-character-basis} również rozkłada charakter iloczynu tensorowego, jeśli znamy charaktery modułów nieprzywiedlnych. Pozostaje znaleźć ogólny wzór na te charaktery, bez uprzedniego konstruowania wszystkich przestrzeni wagowych. Temu służy wzór Weyla.
